\section{Conjugating the reference character and tying it to the core}

% BEGIN C77T1
\paragraph{C77T1. A unitary Hecke character at every required embedding.}
Let $\chi_0$ be the reference row character. Its level pin makes
it trivial both on $E^\times$ and on ideles from the real
subfield $F$. Its finite values have the form $\sigma_0(y)$ for
$y\in K$ by K77F4. For each embedding $\sigma:K\to\mathbb C$
we construct a continuous circle-valued character with finite
values $\sigma(y)$ and the prescribed conjugate infinity type.

Write an idele as $a(x)b(x)$, with $a(x)$ purely archimedean
and $b(x)$ purely finite. On the complex archimedean factors a
continuous character trivial on the real factors descends to a
product of circles $z\mapsto z/\bar z$. Its characters are
integer powers, by lifting a circle character along the real
exponential covering. Thus the archimedean type is a product
of powers $k_\tau$ of these angles. Transport the embeddings of
$E$ in this product by the prescribed relabeling for $\sigma$;
call the resulting continuous unitary character $A_\sigma$.
For $u=e_0 r$ with $e_0\in E^\times$ principal and $r$ from
the real ideles, put
\[
  w=\prod_\tau \operatorname{emb}_{K,\tau}
                  (e_0/\bar e_0)^{k_\tau}.
\]
The real idele contributes angle one. The original and transported
archimedean factors on $u$ are $\sigma_0(w)$ and $\sigma(w)$.
Since $\chi_0(u)=1$, the finite factor equals $\sigma_0(w^{-1})$.

We must show $|\sigma(y)|=1$ when
$\chi_0(b(x))=\sigma_0(y)$. On the product of finite integral
unit groups $\chi_0$ has finite image: its kernel contains an open
subgroup by the small-neighborhood argument of K77F1, and the
unit group is compact. Let $m>0$ kill that image. Raise the
fractional ideal of $b(x)$ to the class number $n$ of $E$, and
write its $n$th power as the principal ideal of $e_0$.
Then $b(x)^n/b(e_0)$ is a finite integral unit. Consequently
$y^{nm}=w^{-m}$ in $K$ by injectivity of $\sigma_0$, with $w$
as above for $e_0$. The absolute value of $\sigma(w)$ is one
since it is an archimedean angle product. This gives
$|\sigma(y)|^{nm}=1$ and hence the assertion.

Now define $\psi_\sigma(x)=\sigma(y)$ using the unique such $y$.
Multiplicativity follows by injectivity of $\sigma_0$.
This character is one on a finite congruence subgroup at the
identity and depends only on the finite part, hence is continuous.
The product $A_\sigma\psi_\sigma$ has the required values on
finite ideles. The calculation with $e_0r$ proves it is one on
principal and on real ideles. Its infinity exponents are precisely
the old exponents in the new ordering. At $\sigma_0$ use the
original $\chi_0$ itself. Carry out this choice for each occupied
row and embedding and use the trivial character at unused indices.
This gives the source family $\chi_W$ with its conjugation law
as well as its prescribed distinguished value.
% END C77T1

% BEGIN C77T2
\paragraph{C77T2. The effect of changing the character and the frame.}
Use the original core with its $\mathrm{LawsJ}$, $\mathrm{LawsCore}$
and algebra pin, the matching base primes and $\mathrm{VolTieG}$,
and use the analytic record's level law and $\mathrm{FrameKs}$.
Take $\chi_W$ from C77T1. At a relevant row and good embedding,
use its determinant-root multiplier $\mu$: for a determinant
$d/\bar d$ use the prescribed value of $\chi_W(d)$ in the source
convention. Two roots differ by a real idele, on which C77T1 gives
one. The resulting unit multiplier is multiplicative; on rational
movers the principal and real character laws give the required
triviality. Use the continuous determinant multiplier construction
of the original core at these embeddings; elsewhere use one.
For the same dictionary the ratio of the new and old multipliers
is measurable and unitary. The squared norm of every adelic cusp
reading is unchanged pointwise, and integration gives the same
Petersson norm and preserves the $L^2$ and definiteness laws.
All valuation tests and algebraic cusp spaces are the original ones.

Use the marked change of P77R1 to install the new base frame and
dictionary. Its right multiplication sends normalized frames to
normalized frames, so the unit-frame law persists. The rational
movers and the level subgroups just identified give the new
adelic automorphy and measurability. For the remaining norm
definiteness use the source volume comparison on these domains:
the adelic measure dominates the component measures on the charts.
Thus zero $L^2$ norm forces every continuous cusp reading to
vanish on its domain, using positive measure on nonempty open
sets; the component supremum definiteness gives zero section.
Square integrability follows from the bounded cusp readings and
finite volume. These are the analytic domain and measure
comparisons used from the lawful-core construction. At the new
dictionary entries the marked lattice equations and integral
sections are as in P77R1. Hence with those comparisons the
modified core retains its two sets of core laws.

In the original model the extremal vector is chosen in the
smooth lowest type. Smoothness is supplied by the source
lowest-type result, using the conjugate model at a reversed
archimedean labeling. Take $\kappa_{\rm Top}$, this extremal
vector, and the marked frame just constructed. Together with the
algebra pin, the core laws and C77T1 these give all six clauses
of $\mathrm{TieIn}$ on the modified core. Its parameter record,
small field, uniformizations and ordered primes are still those
of the starting core. This is the adapted input to P77R2.
% END C77T2
